\section*{Речник на основните термини}
\addcontentsline{toc}{section}{Речник на основните термини}
\label{sec:rechnik}

{\small
\begin{longtable}{@{}T{4.1cm}L{12.7cm}@{}}
Прокурор & Магистрат от Прокуратурата на Република България, който повдига и поддържа
обвинението за престъпления от общ характер и ръководи разследването (чл.~127 КРБ;
чл.~46 НПК).\\
Наблюдаващ прокурор & Прокурорът, който ръководи конкретното разследване и осъществява
постоянен надзор за законосъобразното и своевременното му провеждане (чл.~46, ал.~2, т.~1
и чл.~196, ал.~2 НПК).\\
Досъдебно производство & Първата фаза на наказателния процес по дела за престъпления от
общ характер, в която се събират доказателствата за престъплението и за извършителя му.
Негови органи са прокурорът и разследващите органи (чл.~193 НПК). Делата, които се
образуват по тъжба на пострадалия до съда, нямат досъдебна фаза.\\
Разследващи органи & Следователите, разследващите полицаи, разследващите митнически
инспектори и разследващите инспектори от Комисията за противодействие на корупцията
(чл.~52, ал.~1 НПК).\\
Предварителна проверка & Проверка, която прокурорът извършва лично или възлага на други
органи, за да прецени дали има основание за образуване на досъдебно производство
(чл.~145, ал.~1, т.~2 и 3 ЗСВ).\\
Законен повод & Изчерпателно изброените в чл.~208 НПК обстоятелства, които дават
основание за започване на разследване.\\
Достатъчно данни & Данни, от които може да се направи основателно предположение, че е
извършено престъпление (чл.~211, ал.~1 НПК).\\
Мярка за неотклонение & Мярка за процесуална принуда спрямо обвиняемия: подписка,
гаранция, домашен арест или задържане под стража (чл.~58 НПК).\\
Обвинителен акт & Актът, с който прокурорът повдига обвинение пред съда (чл.~246 НПК).\\
Тежко престъпление & Престъпление, за което законът предвижда лишаване от свобода повече
от пет години, доживотен затвор или доживотен затвор без замяна (чл.~93, т.~7 НК).\\
Информационна система & Всяко отделно устройство или съвкупност от взаимосвързани или
сходни устройства, което осигурява автоматична обработка на данни по определена програма,
както и компютърните данни, които то обработва (чл.~93, т.~21 НК). В НПК се използва
терминът „компютърна информационна система“.\\
Компютърни данни & Всяко представяне на факти, информация или понятия във форма,
поддаваща се на обработка в информационни системи, включително програма (чл.~93, т.~22
НК). В НПК се използва терминът „компютърни информационни данни“.\\
Данни за трафика & Данните, необходими за проследяване и идентифициране на източника и
направлението на връзката, на датата, часа, продължителността и типа ѝ, на крайното
устройство и на използваните клетки (чл.~159а, ал.~1 НПК; чл.~251б, ал.~1 ЗЕС).\\
Електронни доказателства & Данни за абонати, данни за трафика или данни за съдържанието,
съхранявани от доставчик на услуги или от негово име в електронна форма (чл.~3, т.~8 от
Регламент (ЕС) 2023/1543). В по-широк смисъл~-- всяка информация с доказателствено
значение, която се съхранява или предава в цифрова форма.\\
Специални разузнавателни средства & Технически средства и оперативни способи за
документиране на дейността на контролирани лица и обекти, като наблюдение, подслушване,
проследяване и проверка на компютърна информация (чл.~172, ал.~1 НПК).\\
\textit{Dominus litis} & Лат. „господар на делото“; израз, с който се описва ръководното
положение на прокурора в досъдебното производство.\\
\end{longtable}
}
% longtable увеличава брояча на таблиците и без надпис; връщаме го, за да започне номерацията от 1.
\addtocounter{table}{-1}
